\documentclass[12pt]{article}
\usepackage[utf8]{inputenc}
\usepackage[T1]{fontenc}
\usepackage{hyperref}
\title{Sumbwa Language Resources for K-12 Curriculum Development}
\author{Compiled by Claude Sonnet 5.5 (claude-sonnet-5-5) for Andr\'as Kornai}
\date{}
\begin{document}
\maketitle

\section{Introduction}
Sumbwa (Kisumbwa) is a Bantu language of north-western Tanzania, spoken mainly in the Bukombe, Mbogwe and Geita districts, with more speakers in Kahama, Biharamulo and Urambo. Wikipedia gives about 361,000 speakers (2009 Tanzania Language Atlas), while Joshua Project gives much higher figures, so the estimates conflict. Wikipedia describes the language as largely undescribed. No school materials were found.

\section{Primary Educational Resources}
No primers or textbooks in Sumbwa were found.

\section{Linguistic and Reference Materials}
\begin{itemize}
\item Sumbwa language (Wikipedia), which cites the trilingual dictionary by Kulikoyela Kahigi (2008, Sumbwa--English--Swahili, ISBN 9987-691-13-7) and an 1898 grammar: \url{https://en.wikipedia.org/wiki/Sumbwa_language}
\item Sumbwa people (Wikipedia): \url{https://en.wikipedia.org/wiki/Sumbwa_people}
\end{itemize}

\section{Religious and Community Resources}
\begin{itemize}
\item Lutheran Bible Translators: Sumbwa Bible translation project (team selected, workshop held; no completed scripture reported): \url{https://lbt.org/?p=5539}
\item Joshua Project language profile (translation status): \url{https://joshuaproject.net/languages/suw}
\end{itemize}

\section{Digital and Audio Resources}
\begin{itemize}
\item Global Recordings Network, Sumbwa audio recordings (evangelism and Bible teaching, not scripture translation): \url{https://globalrecordings.net/en/language/suw}
\end{itemize}

\section{Curriculum Development Context}
The Kahigi dictionary appears to be the main lexical resource. Orthography, reading materials and teaching grammars were not found. A curriculum would need new materials developed with Sumbwa speakers. The Bible translation effort may eventually yield written text. I did not verify the availability of the dictionary online.

\end{document}